% ── 다리 초안: vanderven1998 → 우리 eq:lco-gxi / eq:lco-spinodal (§13 order-disorder) ──
% 삽입 위치: _sections/ch1_sec13_lcohys.tex, §13.3(T2·T3 order--disorder) 계보 문단의
%   전개 수식 블록(현행 115–123행) 직후, "★$\Omega$ 와 config $\Delta S$ 의 구분(혼동 가드)"
%   (현행 124행) 앞.
% 앵커 문장: "...제일원리 계산이 $x{=}\tfrac12$ 부근 Li$\cdot$빈자리 질서상의 안정 창과 전이를
%   재현했으며\cite{vanderven1998}, ..."(현행 108–110행) 및 "유효 인력 $\Omega_j^\mathrm{cat}>0$
%   (미시 기원은 Li$\cdot$빈자리의 정렬 상호작용...)"(현행 112–113행).
% 컨테이너: srcbox. 우리 기호로 서술, 논문 량은 표에서만.
% ★원문 검증: 방법(제일원리 클러스터 전개 + 상평형)·$x{=}\tfrac12$ Li 질서상·저-$x$ staging
%   = WebSearch 초록·ADS(1998PhRvB..58.2975V) 대조 확인. 정확한 클러스터 전개 에너지식·ECI 값·
%   식번호 = 원문 PDF(구 APS 포맷) 기계 판독 실패 → ★확인 필요(논문 식 재현 안 함, 방법·결과
%   구조 수준 대응). 아래 등치식은 \emph{우리 평균장 축약}을 적은 것이지 원 논문 식이 아니다.

\begin{srcbox}
\textbf{다리 --- 유효 인력 $\Omega_j^\mathrm{cat}$ 가 어느 제일원리 상도표에 대응하는가.} T2$\cdot$T3 를
정규용액~\eqref{eq:lco-gxi} 의 \emph{유효 인력} $\Omega_j^\mathrm{cat}>0$ 로 접을 때, 그 미시 기원인
``Li$\cdot$빈자리의 정렬 상호작용''과 $x{=}\tfrac12$ 부근 질서상의 안정성은 Van der Ven
등\cite{vanderven1998} 의 제일원리 상 안정성 계산이 준다. 대응은 다음과 같다(왼쪽이 본문 기호,
오른쪽이 원 논문 량):
\begin{center}\footnotesize
\begin{tabular}{@{}ll@{}}
\hline
본문(식~\eqref{eq:lco-gxi},~\eqref{eq:lco-spinodal}) & Van der Ven\cite{vanderven1998} 대응량 \\
\hline
유효 인력 $\Omega_j^\mathrm{cat}$ (평균장 쌍) & 유효 클러스터 상호작용(ECI) 집합 $\{V_i,V_{ik},\dots\}$ \\
정규용액 볼록성 상실 $\Omega_j^\mathrm{cat}\!>\!2RT$ & 제일원리 자유에너지의 2상 안정역(상분리) \\
$x{\approx}\tfrac12$ 질서상(T2/T3 monoclinic) & 계산이 재현한 $x{=}\tfrac12$ Li$\cdot$빈자리 질서상 \\
저농도 영역(하프셀 상한 밖 옵션) & 저-$x$ staging 화합물(교대 Li 평면 점유) \\
\hline
\end{tabular}
\end{center}
등치의 핵은 \emph{평균장 축약}이다 --- 원 논문의 다체 클러스터 전개를 전이 $j$ 창의 진행률 $\xi_j$
좌표에서 최근접 쌍 한 항으로 접으면, 그 유효 쌍상호작용이 본문의 $\Omega_j^\mathrm{cat}$ 이다:
\begin{equation}
\underbrace{\Omega_j^\mathrm{cat}\,\xi_j(1-\xi_j)}_{\text{식~\eqref{eq:lco-gxi} 정규용액 항}}
\;\longleftrightarrow\;
\underbrace{\big[\text{ECI 전개의 최근접 쌍 몫}\big]_{\text{평균장}}}_{\text{Van der Ven 클러스터 전개의 축약}},
\qquad
\Omega_j^\mathrm{cat}>2RT\ \Leftrightarrow\ \text{$x{\approx}\tfrac12$ 질서상 안정},
\label{eq:br-vanderven1998-1}
\end{equation}
곧 우리 이중웰 문턱 $\Omega_j^\mathrm{cat}>2RT$ 가 켜지는 조건이, 그들의 상도표가 $x{=}\tfrac12$ 질서상을
안정화하는 조성$\cdot$온도 창에 대응한다(통계역학--연속체 스케일 브리징\cite{ml2024} 이 같은 전이를 연속
조성축에서 재확인). \emph{방법 요지} --- Van der Ven 등은 $\mathrm{Li}_x\mathrm{CoO_2}$ 형성에너지를
제일원리 클러스터 전개(자리 점유 변수의 유효 클러스터 상호작용 합)로 적고 이를 상평형(자유에너지
최소)으로 풀어 $x{=}\tfrac12$ Li$\cdot$빈자리 질서상과 저-$x$ staging 을 재현한 상도표를 얻는다.
\emph{가정 차} --- 원 계산은 여러 ECI 와 \emph{명시적} 질서 바닥상태를 전부 담으나, 본문은 이를 전이별
\emph{단일 평균장 쌍상호작용} $\Omega_j^\mathrm{cat}$ 하나로 축약(정규용액 이중웰)하므로 질서상의 미시
구조가 아니라 상분리의 유효 문턱만 남기고, 각 전이의 실제 gap 유무는 최종적으로 \emph{피팅된}
$\Omega_j^\mathrm{cat}$ 이 $2RT$ 를 넘는지가 정한다(pure-LCO $\Omega_j^\mathrm{cat}$ 는 초기값,
\S\ref{sec:lco-hys-dope}).
\end{srcbox}
